\documentclass[10pt,a4paper]{amsart}
\usepackage[margin=19mm]{geometry}
\usepackage{amsmath,amssymb,mathtools}
\usepackage[scr=rsfs,cal=euler]{mathalfa}
\usepackage{xfrac}
\usepackage[unicode,hidelinks,bookmarks=false]{hyperref}
\newcommand{\ie}{i.e.\ }
\newcommand{\cf}{cf.\ }
\newcommand{\resp}{respectively\ }
\newcommand{\Subsection}{\subsection}
\providecommand{\nobreakdash}{\nobreak\hbox{-}}
\setlength{\emergencystretch}{2em}
\multlinegap0pt
\usepackage{mathtools}
\usepackage{amssymb}
\usepackage{enumerate}
\newtheorem{lemma}{Lemma}
\newtheorem{corollary}{Corollary}
\newcommand{\Gm}{\mathbb{G}_m}
\newcommand{\K}{\mathbb{K}}
\newcommand{\LL}{\mathbb{L}}
\newcommand{\TT}{\mathbb{T}}
\newcommand{\cO}{\mathcal{O}}
\newcommand{\dArt}{\mathbf{dArt}}
\title{Equivariant Quotients of derived symplectic spaces and \\ Legendrian Intersection Theorem}
\author{Efe \.{I}zbudak \and Kadri \.{I}lker Berktav}
\date{Source: arXiv:2605.08394v2 (CC BY 4.0)}
\begin{document}
\maketitle

\noindent\emph{Source and licence.} Equivariant Quotients of derived symplectic spaces and \\ Legendrian Intersection Theorem. Version arXiv:2605.08394v2 (CC BY 4.0), distributed by arXiv under the Creative Commons Attribution 4.0 International licence, http://creativecommons.org/licenses/by/4.0/, as recorded in the arXiv licence metadata for this exact version and shown on the abstract page https://arxiv.org/abs/2605.08394v2. Full licence notice: \texttt{licenses/arxiv-2605.08394-CC-BY-4.0.txt}.\par\smallskip

\noindent\textit{Editorial scope.} Counted unit: the article's corollary cor:fundamental\_morphism (source0.tex lines 974--979). The article's complete original proof of it is delivered with it (source0.tex lines 980--1001, one \texttt{proof} environment of 22 lines). No mathematics is written, repaired or re-derived: the statement and the proof are the article's own lines, in the article's own notation, and no retained line is substituted or re-flowed.  Retained uncounted as the source-local context the proof needs: lines 906--911.  Nothing was omitted from any retained span, so the package carries no omission record.  No retained line contains a citation, so the wrapper carries no bibliography and no bibliography entry.  No retained line contains a figure environment, an image-inclusion command or drawing code, so no image is delivered and the package declares no asset dependency.  Editorial formatting only, in addition: an AMS \texttt{amsart} preamble that restates the article's own declarations and macros used by the retained lines, namely the environments \texttt{lemma}, \texttt{corollary} on separate counters, together with the standard packages those lines need.  The retained environments print in this document as Lemma 1, Corollary 1 in order of appearance, whereas the article prints the same items with the numbering of its own full document.  The complete native source of the article as distributed in the arXiv e-print for this exact version is bundled unchanged as \texttt{sources/200150/source0.tex}.
\begin{lemma}[Derived Atiyah Sequence] \label{lem:atiyah}
  Let $G$ be a smooth affine group scheme with Lie algebra $\mathfrak{g}$. For a principal $G$-torsor $p \colon \widetilde{X} \to X$ in $\dArt_\K$, the $G$-action induces a stable fiber sequence in $L_{qcoh}(\widetilde{X})$
  \begin{equation*}
    \mathbb{L}p^*\LL_{X} \to \LL_{\widetilde{X}} \to \mathfrak{g}^\vee \otimes_\K \cO_{\widetilde{X}}.
  \end{equation*}
\end{lemma}

\begin{corollary} \label{cor:fundamental_morphism}
  Let $p \colon \widetilde{X} \to X$ be a principal $\Gm$-torsor in $\dArt_\K$. The standard left-invariant vector field on $\Gm$ determines a fundamental morphism $Y \colon \cO_{\widetilde{X}} \to \TT_{\widetilde{X}}$ in $L_{qcoh}(\widetilde{X})$ such that the relative cotangent sequence takes the form
  \begin{equation*}
    \LL p^* \LL_X \xrightarrow{(\TT p)^\vee} \LL_{\widetilde{X}} \xrightarrow{Y^\vee} \cO_{\widetilde{X}}.
  \end{equation*}
\end{corollary}

\begin{proof}
  The proof of Lemma \ref{lem:atiyah} establishes the equivalence $\LL_{\widetilde{X}/X} \simeq \mathfrak{g}^\vee \otimes_\K \cO_{\widetilde{X}}$. Taking the derived $\cO_{\widetilde{X}}$-linear dual of this equivalence gives an isomorphism
  \begin{equation*}
    \TT_{\widetilde{X}/X} \simeq \mathfrak{g} \otimes_\K \cO_{\widetilde{X}}.
  \end{equation*}
  The relative tangent complex sits in the transitivity fiber sequence $\TT_{\widetilde{X}/X} \to \TT_{\widetilde{X}} \xrightarrow{\TT p} \LL p^* \TT_X$. Composing the isomorphism with the canonical morphism $\TT_{\widetilde{X}/X} \to \TT_{\widetilde{X}}$ defines the infinitesimal action morphism
  \begin{equation*}
    a \colon \mathfrak{g} \otimes_\K \cO_{\widetilde{X}} \to \TT_{\widetilde{X}}.
  \end{equation*}
  The multiplicative group scheme $\Gm$ has the coordinate algebra $\K[t, t^{-1}]$. The tangent space at the identity is the $1$-dimensional $\K$-vector space generated by the derivation $\partial = t \partial_t$. This generator defines an isomorphism of $\K$-vector spaces $\K \xrightarrow{\sim} \mathfrak{g}$. The fundamental morphism $Y$ is defined as the composition
  \begin{equation*}
    Y \colon \cO_{\widetilde{X}} \simeq \K \otimes_\K \cO_{\widetilde{X}} \xrightarrow{\sim} \mathfrak{g} \otimes_\K \cO_{\widetilde{X}} \xrightarrow{a} \TT_{\widetilde{X}}.
  \end{equation*}
  Taking the derived dual of this composition we obtain
  \begin{equation*}
    Y^\vee \colon \LL_{\widetilde{X}} \xrightarrow{a^\vee} \mathfrak{g}^\vee \otimes_\K \cO_{\widetilde{X}} \xrightarrow{\sim} \cO_{\widetilde{X}}.
  \end{equation*}
  The morphism $a^\vee$ is the canonical projection $\LL_{\widetilde{X}} \to \LL_{\widetilde{X}/X}$ in the cotangent transitivity sequence, post-composed with the equivalence from Lemma \ref{lem:atiyah}. Substituting this identification into the standard transitivity sequence gives the stable fiber sequence
  \begin{equation*}
    \LL p^* \LL_X \xrightarrow{(\TT p)^\vee} \LL_{\widetilde{X}} \xrightarrow{Y^\vee} \cO_{\widetilde{X}}.
  \end{equation*}
\end{proof}

\end{document}
